\documentclass[../main.tex]{subfiles}
\graphicspath{{\subfix{../}}}
\begin{document}

\textit{Virtuális munka elve, Hamilton-elv. Legkisebb hatás elve. Lagrange-féle elsőfajú és másodfajú mozgásegyenletek. Hamilton-függvény, kanonikus egyenletek. Kanonikus transzformációk. Szimmetriák és megmaradási tételek.}

\subsection{Bevezető}
A klasszikus mechanika alapját a Newtoni formalizumus képzi, ami a Newton törvények és megmaradási törvények segítségével írja le a testek mozgását. Ennek alapját képzi a Newton törvények, amelyek a következők:
\begin{itemize}
    \item Inerciarendszer: Minden test, amelyre nem hat külső erő, egyenes vonalú egyenletes mozgást végez.
    \item Erő és gyorsulás: A testre ható erő egyenlő a test tömegének és gyorsulásának szorzatával: $\bold{F} = m \bold{a}$.
    \item Akció és reakció: Minden erőhatásra van egy vele ellentétes irányú és egyenlő nagyságú erőhatás.
\end{itemize}
És ezekből levezethetők a megmaradási törvények:
\begin{itemize}
    \item Energia megmaradás: Egy zárt rendszerben az összenergia állandó $E = K + U = const$, ahol $K$ a kinetikus energia és $U$ a potenciális energia.
    \item Impulzus megmaradás: Egy zárt rendszerben az összimpulzus állandó $\bold{P} = \sum m_i \bold{v}_i = const$.
    \item Forgásimpulzus megmaradás: Egy zárt rendszerben az összforgásimpulzus állandó $\bold{L} = \sum \bold{r}_i \times m_i \bold{v}_i = const$.
\end{itemize}
Ezeket a törvényeket felhasználva minden mechanikai problémára fel lehet írni a mozgásegyenleteket, amelyet egy differenciálegyenlet segítségével megoldhatunk.

Ez viszont gyakran olyan egyenleteket erredményez, amelyeket nagyon nehéz vagy analitikusan akár lehetetlen megoldani. Ezért az évszázadok során megpróbálták ezeket a fundamentális törvényeket más szempontból megközelíteni,
és új fomalizmusokat létrehozni. Így járt el például Lagrange és Hamilton is a 18-19. században, akik a mechanikát egy új megközelítésből írták le, amely a mozgásegyenletek helyett a rendszer energiájára és impulzusára fókuszált.
Ők abból az irányból közelítették meg a problémát, hogy a fizikának látszólag van egy olyan hajlama, hogy egy rendszer egy adott állapotból egy másikba úgy menjen át, hogy mindig a "legrövidebb" utat válassza, a két pont közötti
végtelen lehetőség közül. Ezt a "legrövidebb" utat pedig úgy definiálták, hogy a rendszer egy adott állapotból egy másikba úgy menjen át, hogy a hatás (action) nevű mennyiség minimális legyen. Ezt a megközelítést Hamilton vezette be, és ezért nevezzük Hamilton-elvnek.
Ebből a megközelítésből kiindulva Lagrange is kidolgozta a saját formalizmusát, amely a Lagrange-féle mozgásegyenleteket eredményezte. Ezek az egyenletek a kinetikus és potenciális energia különbségéből származnak, és lehetővé teszik a
mozgásegyenletek felírását anélkül, hogy közvetlenül az erőket kellene figyelembe venni. Ez különösen hasznos olyan rendszerek esetén, ahol az erők bonyolultak vagy nehezen meghatározhatók.
A Hamilton és Lagrange formalizmusok tehát alternatív megközelítést kínálnak a klasszikus mechanika problémáinak megoldására, és gyakran egyszerűbbé teszik a mozgásegyenletek felírását és megoldását.

\subsection{Virtuális munka elve, Hamilton-elv. Legkisebb hatás elve.}
\subsubsection{Virtuális munka elve}
A virtuális munka elve azon a megfigyelésen alapszik, hogy ha egy testre erő hat, és emiatt $A$ pontból $B$ pontba mozog, akkor a két pont között úgy választja ki a pályát, hogy a pálya és a körülötte lévő "közeli" pályák között a a munka nulla legyen (elsőrendűen).
Ez tehát a munka függvényének egy szélsőértékét kell megtalálni  (lehet inflexiós pont is) ami azt jelenti, hogy kis változásokra munka nem változik elsőrendűen. A munkát úgy definiáltuk, hogy:
\begin{equation}
    W = \int_{\bold{r}(t_0) = A}^{\bold{r}(t_1) = B} \bold{F} \cdot d\bold{r} = \int_{t_0}^{t_1} \bold{F} \cdot \frac{d\bold{r}}{dt} dt = \int_{t_0}^{t_1} \bold{F} \cdot \bold{v} dt
\end{equation}
ahol $\bold{F}$ a testre ható erő, és $d\bold{r}$ a test kis elmozdulása. Most tegyük fel, hogy a test egy kicsit eltér a valódi pályától, és egy közeli pályán mozog, amelyet $\bold{r}(t) + \delta \bold{r}(t)$-vel jelölünk.
Ekkor a munka változása a két pálya között a következőképpen írható fel:
\begin{equation}
    \delta W = \int_{t_0}^{t_1} \bold{F} \cdot \frac{d}{dt} (\bold{r} + \delta \bold{r}) dt - \int_{t_0}^{t_1} \bold{F} \cdot \frac{d\bold{r}}{dt} dt
\end{equation}
\begin{equation}
    \delta W = \int_{t_0}^{t_1} \bold{F} \cdot \frac{d (\delta \bold{r})}{dt} dt
\end{equation}
Ezt nevezik virtuális munkának, és ez a munka egy olyan kis virtuális elmozdulásra vonatkozik, amely nem valósul meg a tényleges rendszeben de elképzelhető.

Ha a rendszer statikus egyensúlyi állapotban van, akkor a virtuális munka elve szerint a virtuális munka nulla kell legyen minden lehetséges virtuális elmozdulásra:
\begin{equation}
    \delta W = \sum_i \bold{F}_i \cdot \delta \bold{r}_i = 0
\end{equation}
ahol $\delta W$ a virtuális munka, $\bold{F}_i$ a rendszerre ható erők, és $\delta \bold{r}_i$ a testek virtuális elmozdulásai. 

Szabad mozgás esetén a feltétel csupán annyi, hogy minden i-re igaz legyen, hogy $\bold{F}_i = 0$. Ha a testeknek viszont vannak kényszerei (pl. egy kötött pályán mozognak), akkor a kényszererők is hatnak a testekre,
és ezeket is figyelembe kell venni a virtuális munka elvében. Tegyük fel, hogy a kényszereket megadhatjuk a következő alakban:
\begin{equation}
    \phi_j(\bold{r}_1, \bold{r}_2, \ldots, \bold{r}_N) = 0, \quad j = 1, 2, \ldots, k
\end{equation}
Ahol $k$ a kényszerek száma. Ezeknek a kényszereknek kis elmozdulás esetén is teljesülniük kell, tehát a kényszerfeltételek teljesülnek a virtuális elmozdulásokra is:
\begin{equation}
    \phi_j (\bold{r}_1 + \delta \bold{r}_1, \bold{r}_2 + \delta \bold{r}_2, \ldots, \bold{r}_N + \delta \bold{r}_N) = 0
\end{equation}
Vezessünk be egy egyszerűbb jelölést a kényszerfeltételekhez:
\begin{equation}
    \{\bold{r}_1, \bold{r}_2, \ldots, \bold{r}_N\} = \bold{r}_i
\end{equation}
\begin{equation}
    \{\bold{r}_1 + \delta \bold{r}_1, \bold{r}_2 + \delta \bold{r}_2, \ldots, \bold{r}_N + \delta \bold{r}_N\} = \bold{r}_i + \delta \bold{r}_i
\end{equation}
Ekkor a kényszerfeltételek a következőképpen írhatók fel:
\begin{equation}
    \phi_j (\bold{r}_i) = 0
\end{equation}
\begin{equation}
    \phi_j (\bold{r}_i + \delta \bold{r}_i) = 0
\end{equation}
Most alkalmazzuk a Taylor-sorfejtést a kényszerfeltételekre, és hagyjuk el a másod- és magasabb rendű tagokat, mivel a virtuális elmozdulások kicsik:
\begin{equation}
    \phi_j (\bold{r}_i + \delta \bold{r}_i) = \phi_j (\bold{r}_i) + \sum_i \frac{\partial \phi_j}{\partial \bold{r}_i} \cdot \delta \bold{r}_i = 0
\end{equation}
Mivel $\phi_j (\bold{r}_i) = 0$, ezért a fenti egyenlet egyszerűsödik:
\begin{equation}
    \sum_i \frac{\partial \phi_j}{\partial \bold{r}_i} \cdot \delta \bold{r}_i = 0
\end{equation}
Vagy a Nabla jelöléssel:
\begin{equation}
    \sum_i \boldsymbol{\nabla}_i \phi_j(\bold{r}_i) \cdot \delta \bold{r}_i = 0
\end{equation}
Ezt az egyenletet a Lagrange multiplikátor módszerével oldhatjuk meg, amelynek használatával megtaláhatjuk egy függvény szélsőértékét kényszerfeltételek mellett. A módszer lényege, hogy bevezetünk egy új függvényt, amely a keresett függvény és a kényszerfeltételek lineáris kombinációja:
\begin{equation}
    \mathcal{L} = \delta W + \sum_j \lambda_j \sum_i \boldsymbol{\nabla}_i \phi_j(\bold{r}_i) \cdot \delta \bold{r}_i
\end{equation}
ahol $\lambda_j$ a Lagrange multiplikátorok. Most már alkalmazhatjuk a virtuális munka elvét a $\mathcal{L}$ függvényre:
\begin{equation}
    \mathcal{L} = \sum_i \bold{F}_i \cdot \delta \bold{r}_i + \sum_j \lambda_j \sum_i \boldsymbol{\nabla}_i \phi_j(\bold{r}_i) \cdot \delta \bold{r}_i = 0
\end{equation}
Mivel a virtuális elmozdulások tetszőlegesek, ezért a fenti egyenlet csak akkor teljesülhet, ha az egyes tagok külön-külön is nulla:
\begin{equation}
    \bold{F}_i + \sum_j \lambda_j \boldsymbol{\nabla}_i \phi_j(\bold{r}_i) = 0
\end{equation}
Ez az egyenlet a rendszer mozgásegyenleteit írja le kényszerfeltételek mellett, ahol a $\lambda_j$ multiplikátorok a kényszererőket reprezentálják.

\subsubsection{Hamilton-elv, Legkisebb hatás elve}
Amikor egy fizikai rendszer változását vizsgáljuk, gyakran előfordul, álapotok között, gyakran több lehetséges útvonal is létezik. Például, ha egy Konzervatív erőtérben (pl. gravitációs térben) leejtünk egy tárgyat,
Akkor munkát a tér csak a vertikális tengely mentén végez, horizontális elmozdulás során a munkavégzés nulla. Tehát elméletileg a tárgy végtelen útvonal közül bármelyiken eljuthat a kiinduló pontból a végpontba, ha csak a vertikális komponensét nézzük.
Felmerülhet tehát a kérdés, hogy ha kísérletet teszünk a rendszer viselkedésének megértésére, és a kezdeti paramétereket mindig azonosnak választjuk, akkor miért mindig ugyanazt az útvonalat követi a rendszer? Miért nem választ véletlenszerűen egy másik lehetséges útvonalat?

A kérdés megválaszolását először a fény útvanálanak vizsgálata indukálta, ahol Fermat felismerte, hogy a fény mindig azt az útvonalat választja, amelyen a fénynek a legkisebb idő alatt kell megtennie a távolságot. Ezt az elvet Fermat elvének nevezzük.
Itt felmerülhet a kérdés, hogy mit is jelent az, hogy "legrövidebb útvonal"? 

Arra hogy egy tetszőleges koordináta rendszerben meg tudjuk határozni a két pont közötti utat, a variációs kalkulust használjuk. A variációs kalkulus egy olyan matematikai eszköz, amely lehetővé teszi a függvények szélsőértékeinek meghatározását, amikor a függvények nem csak egy változótól, hanem egy függvénytől is függenek.
A variációs kalkulus alapvető fogalma a funkcionál, amely egy olyan függvény, amely egy függvényt rendel egy valós számhoz\footnote{
    Kicsit pontosabban a funkcionál egy olyan objektum, ami egy függvényt rendel egy valós számhoz, és általában egy integrál formájában van megadva, ahol a függvény és annak deriváltjai szerepelnek az integrandumban. Például:
    \begin{equation}
        S[f] = \int_{0}^{1} f'(x)^2 \, dx
    \end{equation}
    Az $S[f]$ értéke egy valós szám amit az $f$ függvényhez rendelünk. 

}. Például, ha van egy $y(x)$ függvényünk, akkor egy funkcionál lehet a következő:
\begin{equation}
    S[y] = \int_{x_1}^{x_2} F(x, y, y') \, dx
\end{equation}
ahol $F$ egy adott függvény, $y'$ pedig a $y$ függvény deriváltja. A funkcionál szélsőértékének meghatározásához a variációs kalkulus a következő lépéseket követi:
\begin{itemize}
    \item Válasszunk egy kis változást a függvényben: $y(x) \to y(x) + \delta y(x)$, ahol $\delta y(x)$ egy kis perturbáció.
    \item Számítsuk ki a funkcionál változását a perturbáció hatására: $\delta S = S[y + \delta y] - S[y]$.
    \item Állítsuk be a változást úgy, hogy a funkcionál változása nulla legyen: $\delta S = 0$.
\end{itemize}
\begin{figure}
    \centering
    \includegraphics[width=0.5\textwidth]{imgs/2-tetel/CalculusofVariations.png}
    \caption{Variációs kalkulus}
    \label{fig:variational_calculus}
\end{figure}
Tehát, ha van két pontom $A = (x_1,y_1)$ és $B = (x_2,y_2)$, akkor a két pont között húzott görbe mentén a következő funkcionált definiálhatjuk:
\begin{equation}
    dS[y] = \sqrt{dx^2 + dy^2}
\end{equation}
\begin{equation}
    dS[y] = \sqrt{1 + \left(\frac{dy}{dx}\right)^2} \, dx
\end{equation}
\begin{equation}
    S[y] = \int_{x_1}^{x_2} \sqrt{1 + \left(\frac{dy}{dx}\right)^2} \, dx
\end{equation}
Ez a funkcionál a két pont közötti görbe hosszát méri. Most vezessünk be egy kis perturbációt a görbébe:
\begin{equation}
    y(x) \to y(x) + \delta y(x)
\end{equation}
Akkor a funkcionál változása a következőképpen alakul:
\begin{equation}
    \delta S = S[y + \delta y] - S[y]
\end{equation}
Ez a görbe "variálása", és a görbe hosszának változását méri a perturbáció hatására. Ha a görbe variációja nulla, akkor az azt jelenti, hogy a görbe hosszának változása a kis perturbáció hatására nulla,
vagyis a görbe egy szélsőértéket képvisel. Ez a szélsőérték lehet minimum, maximum vagy egy inflexiós pont is. Tehát, ha a görbe minimumát akarjuk megtalálni, akkor ennek a variált Funkciónálnak a változását nullára kell állítanunk:
\begin{equation}
    \delta S = \delta \int_{x_1}^{x_2} \sqrt{1 + \left(\frac{dy}{dx}\right)^2} \, dx = 0
\end{equation}
Ez a variációs számítás alapja.

\subsubsection*{Példa: Fermat-elv}
Nézzük meg, hogy akkor a Fermat-elv ismeretében hogyan tudjuk meghatározni a fény útvonalát két pont között. Tegyük fel, hogy a fény egy közegben terjed, ahol a fénysebesség $v$ és a közeg törésmutatója $n = \frac{c}{v}$, ahol $c$ a vákuumbeli fénysebesség.
Tudjuk, hogy a fény az idő "szélsőértékét keresi", tehát a két pont között a következő időt kell minimalizálnunk:
\begin{equation}
    \delta t = 0
\end{equation}
Ahol $t$ a fény által megtett idő a két pont között. A fény által megtett idő a következőképpen számítható ki:
\begin{equation}
    t = \int_{A}^{B} \frac{ds}{v} = \int_{A}^{B} \frac{n \, ds}{c}
\end{equation}
Ahol $ds$ a fény által megtett útszakasz. A fény által megtett útszakasz a következőképpen írható fel:
\begin{equation}
    ds = \sqrt{dx^2 + dy^2} = \sqrt{1 + \left(\frac{dy}{dx}\right)^2} \, dx
\end{equation}
Tehát a fény által megtett idő a következőképpen alakul:
\begin{equation}
    t = \int_{x_1}^{x_2} \frac{n}{c} \sqrt{1 + \left(\frac{dy}{dx}\right)^2} \, dx
\end{equation}
Most alkalmazzuk a variációs kalkulus elveit, és vezessünk be egy kis perturbációt a görbébe:
\begin{equation}
    y(x) \to y(x) + \delta y(x)
\end{equation}
Akkor a fény által megtett idő változása a következőképpen alakul:
\begin{equation}
    \delta t = \delta \int_{x_1}^{x_2} \frac{n}{c} \sqrt{1 + \left(\frac{dy}{dx}\right)^2} \, dx = 0
\end{equation}
Ez a variációs számítás alapja a Fermat-elv alkalmazásának a fény útvonalának meghatározására. A variációs számítás elvégzése után megkapjuk a fény útvonalát, amely a két pont között a legkisebb idő alatt teszi meg az utat.

\subsubsection*{Maupertuis-elv}
Pierre Louis Maupertuis, francia matematikus és filozófus, a 18. században dolgozott ezen a poblémán és póbálta meg általánosítani a Fermat-elvet minden fizikai rendszerre. Maupertuis felismerte, hogy a fizikai rendszerek hajlamosak arra,
hogy egy adott állapotból egy másik állapotba úgy menjenek át, hogy egy funkcionál minimális legyen. Ezt a funkcionált Maupertuis hatásnak nevezte el, és a következőképpen definiálta:
\begin{equation}
    \bold{S}[\bold{q}(t)] = \int \bold{p} \, d\bold{q}
\end{equation}
ahol $\bold{p}$ a rendszer általánosított impulzusa (nem feltétlen egyezik meg a fizikai impulzussal), és $\bold{q}$ a rendszer általánosított koordinátái. Maupertuis elve azt mondja ki, hogy egy fizikai rendszer akkor megy át egy adott állapotból egy másik állapotba,
ha a Maupertuis hatás minimális:
\begin{equation}
    \delta \bold{S} = 0
\end{equation}
Ez az elv hasonló a Fermat-elvhez, de általánosabb, mivel nem csak a fényre, hanem minden fizikai rendszerre alkalmazható. Maupertuis elve tehát egy univerzális elv, amely a fizikai rendszerek viselkedését írja le.

\subsubsection*{Hamilton-elv}
A Hamilton-elv hasonló a Maupertuis-elvhez, de egy kicsit általánosabb formában. Hamilton elve azt mondja ki, hogy egy fizikai rendszer akkor megy át egy adott állapotból egy másik állapotba,
ha a Lagrange-függvény integrálja, amelyet akciónak nevezünk, minimális:
\begin{equation}
    S = \int_{t_1}^{t_2} \mathcal{L}(q, \dot{q}, t) \, dt
\end{equation}
Itt az $S$ szintén az akciót jelöli. Ahogy hogy ezt minimalizálni tudjuk, szintén a variációs kalkulus eszközeit használjuk:
\begin{equation}
    \delta S = \delta \int_{t_1}^{t_2} \mathcal{L}(q, \dot{q}, t) \, dt = 0
\end{equation}
A Hamilton-elv tehát egy univerzális elv, amely a fizikai rendszerek viselkedését írja le, és amely a Lagrange-függvényre épül. Ez az elv lehetővé teszi a mozgásegyenletek felírását anélkül, hogy közvetlenül az erőket kellene figyelembe venni.

Ahhoz, hogy meg tudjuk oldani a Hamilton-elv által felírt variációs problémát, bontsuk fel az eseményeket egy kis időintervallumra:
\begin{figure}[H]
    \centering
    \includegraphics[width=0.5\textwidth]{imgs/2-tetel/Hamilton1.png}
    \caption{Diszkretizáció}
    \label{fig:discretization}
\end{figure}
\begin{equation}
    \int dt = \epsilon \sum_i \mathcal{L}(q_i, \frac{q_{i+1} - q_i}{\epsilon})
\end{equation}
Ekkor ha csak egy időintervallumot nézünk, akkor a következő kifejezést kapjuk:
\begin{figure}[H]
    \centering
    \includegraphics[width=0.5\textwidth]{imgs/2-tetel/Hamilton2.png}
    \caption{Diszkretizáció 2}
    \label{fig:discretization2}
\end{figure}
\begin{equation}
    \epsilon\mathcal{L}(q_i, \frac{q_{i+1} - q_i}{\epsilon}) + \epsilon\mathcal{L}(q_{i-1}, \frac{q_{i} - q_{i-1}}{\epsilon})
\end{equation}

Ennek a kifejezésnek vagyunk kiváncsiak, hogyan változik egy kis perturbáció hatására. Ezért deriváljuk a kifejezést $q_i$ szerint:
\begin{equation}
    \frac{\partial}{\partial q_i} \left( \epsilon\mathcal{L}(q_i, \frac{q_{i+1} - q_i}{\epsilon}) + \epsilon\mathcal{L}(q_{i-1}, \frac{q_{i} - q_{i-1}}{\epsilon}) \right) = 0
\end{equation}
Vezessük be a kifejezést $v_i = \frac{q_{i+1} - q_i}{\epsilon}$ sebességet, és ekkor a következő eredményt kapjuk:
\begin{equation}
    \begin{aligned}
        \frac{\partial}{\partial q_i} \epsilon\mathcal{L}(q_i, v_i) + \frac{\partial}{\partial q_i} \epsilon\mathcal{L}(q_{i-1}, v_{i-1}) &= 0 \\
        \epsilon \frac{\partial \mathcal{L}}{\partial q_i} + \epsilon \frac{\partial \mathcal{L}}{\partial v_i} \frac{\partial v_i}{\partial q_i} + \epsilon \frac{\partial \mathcal{L}}{\partial v_{i-1}} \frac{\partial v_{i-1}}{\partial q_i} &= 0
    \end{aligned}
\end{equation}
Itt a többváltozós láncszabályt alkalmaztuk\footnote{
    A láncszabály szerint, ha egy függvény több változótól függ, és ezek a változók maguk is függnek egy harmadik változótól, akkor a teljes derivált kiszámításához minden változó részleges deriváltját ki kell számítani, és ezeket össze kell adni:
\begin{equation}
    \frac{d f(y(x), z(x), \ldots)}{d x} = \frac{\partial f}{\partial y} \frac{d y}{d x} + \frac{\partial f}{\partial z} \frac{d z}{d x} + \ldots
\end{equation}
}
Az általánosított sebességek deriváltjai a következőképpen alakulnak:
\begin{equation}
    \frac{\partial v_i}{\partial q_i} = \frac{\partial}{\partial q_i} \left( \frac{q_{i+1} - q_i}{\epsilon} \right) =  - \frac{1}{\epsilon} \quad \text{és} \quad \frac{\partial v_{i-1}}{\partial q_i} = \frac{\partial}{\partial q_i} \left( \frac{q_{i} - q_{i-1}}{\epsilon} \right) = \frac{1}{\epsilon}
\end{equation}
Vagyis a perturbáció hatására a következő egyenletet kapjuk:
\begin{equation}
    \epsilon \left( \frac{\partial \mathcal{L}}{\partial q_i} - \frac{1}{\epsilon} \frac{\partial \mathcal{L}}{\partial v_i} + \frac{1}{\epsilon} \frac{\partial \mathcal{L}}{\partial v_{i-1}} \right) = 0
\end{equation}
Ahol
\begin{equation}
    - \frac{1}{\epsilon} \frac{\partial \mathcal{L}}{\partial v_i} + \frac{1}{\epsilon} \frac{\partial \mathcal{L}}{\partial v_{i-1}} = \frac{d}{dt} \frac{\partial \mathcal{L}}{\partial v_i}
\end{equation}
Amennyiben $\epsilon$-t infinitezimálisan kicsivé tesszük. Tehát a Hamilton-elvből a következő egyenletet kapjuk:
\begin{equation}
    \frac{\partial \mathcal{L}}{\partial q_i} - \frac{d}{dt} \frac{\partial \mathcal{L}}{\partial v_i} = 0
\end{equation}
vagy a klasszikus jelöléssel:
\begin{equation}
    \frac{\partial \mathcal{L}}{\partial q} - \frac{d}{dt} \frac{\partial \mathcal{L}}{\partial \dot{q}} = 0
\end{equation}
Ez az Euler-Lagrange-féle mozgásegyenlet, amely a Hamilton-elvből származtatható. Ez az egyenlet lehetővé teszi a mozgásegyenletek felírását anélkül, hogy közvetlenül az erőket kellene figyelembe venni. Az egyenletben megadott koordináták és sebességek általánosítottak,
tehát nem feltétlenül egyeznek meg a fizikai koordinátákkal és sebességekkel.

Nézzük még meg, hogy a Hamilton-elvből hogyan vezethető le a Newton-féle mozgásegyenlet. Tegyük fel, hogy egy testre hat egy konzervatív erőtér, amelyet egy potenciális energiafüggvény ír le:
\begin{equation}
    U = U(q)
\end{equation}
Akkor a Lagrange-függvény a következőképpen írható fel:
\begin{equation}
    \mathcal{L} = K - U = \frac{1}{2}m \dot{q}^2 - U(q)
\end{equation}
Ahol $K$ a kinetikus energia, $U$ a potenciális energia, $q$ pedig a koordináta. Most alkalmazzuk az Euler-Lagrange-féle mozgásegyenletet:
\begin{equation}
    \frac{\partial \mathcal{L}}{\partial q} - \frac{d}{dt} \frac{\partial \mathcal{L}}{\partial \dot{q}} = 0
\end{equation}
Számítsuk ki az egyes tagokat:
\begin{equation}
    \frac{\partial \mathcal{L}}{\partial q} = - \frac{\partial U}{\partial q}
\end{equation}
\begin{equation}
    \frac{\partial \mathcal{L}}{\partial \dot{q}} = m \dot{q}
\end{equation}
Ahol a kinetikus energia deriváltja a sebesség szerint a kanonikus impulzus, amely ebben az esetben $p = m \dot{q}$. Most vegyük ennek a deriváltját az idő szerint:
\begin{equation}
    \frac{d}{dt} \frac{\partial \mathcal{L}}{\partial \dot{q}} = m \ddot{q}
\end{equation}
Most helyettesítsük be ezeket az egyenletbe:
\begin{equation}
    - \frac{\partial U}{\partial q} - m \ddot{q} = 0
\end{equation}
Ezt átrendezve megkapjuk a Newton-féle mozgásegyenletet:
\begin{equation}
    m \ddot{q} = - \frac{\partial U}{\partial q} = \bold{F}
\end{equation}
Ez az egyenlet azt mondja ki, hogy a test gyorsulása arányos a rá ható erővel, amelyet a potenciális energia gradiensével adhatunk meg. Tehát a Hamilton-elvből kiindulva sikerült levezetnünk a klasszikus mechanika alapegyenletét, a Newton-féle mozgásegyenletet.

\subsubsection*{Példa: Forgó koordináta rendszer}
Vegyünk egy koodináta rendszert, amely egy tengely körül forog egy szögsebességgel $\omega$. Ebben a rendszerbben a koordináták:
\begin{equation}
    x' = x \cos(\omega t) + y \sin(\omega t)
\end{equation}
\begin{equation}
    y' = -x \sin(\omega t) + y \cos(\omega t)
\end{equation}
Ekkor a Lagrange-függvény a mozgó koordinátákra a következőképpen alakul:
\begin{equation}
    \mathcal{L} = \frac{1}{2}m(\dot{x'}^2 + \dot{y'}^2) - U(x', y')
\end{equation}
Számoljuk ki a derriváltjait a gyorsuló koordinátákra:
\begin{equation}
    \dot{x'} = \dot{x} \cos(\omega t) + \dot{y} \sin(\omega t) - \omega x \sin(\omega t) + \omega y \cos(\omega t)
\end{equation}
\begin{equation}
    \dot{y'} = -\dot{x} \sin(\omega t) + \dot{y} \cos(\omega t) - \omega x \cos(\omega t) - \omega y \sin(\omega t)
\end{equation}
Most helyettesítsük be ezeket:
\begin{equation}
    \begin{aligned}
        \dot{x'}^2 + \dot{y'}^2 = & (\dot{x} \cos(\omega t) + \dot{y} \sin(\omega t) - \omega x \sin(\omega t) + \omega y \cos(\omega t))^2 + \\
        & (-\dot{x} \sin(\omega t) + \dot{y} \cos(\omega t) - \omega x \cos(\omega t) - \omega y \sin(\omega t))^2
    \end{aligned}
\end{equation}
Ezt kibontva és egyszerűsítve a következő eredményt kapjuk:
\begin{equation}
    \dot{x'}^2 + \dot{y'}^2 = \dot{x}^2 + \dot{y}^2 + \omega^2 (x^2 + y^2) + 2\omega (x \dot{y} - y \dot{x})
\end{equation}
Most helyettesítsük be ezt a Lagrange-függvénybe:
\begin{equation}
    \mathcal{L} = \frac{m}{2}(\dot{x}^2 + \dot{y}^2) + \frac{m\omega^2}{2}(x^2 + y^2) + \frac{m\omega}{2}(x \dot{y} - y \dot{x}) - U(x, y)
\end{equation}
Ahol $\frac{m\omega^2}{2}(x^2 + y^2)$ a centrifugális erő, és $\frac{m\omega}{2}(x \dot{y} - y \dot{x})$ a Coriolis erő. Most alkalmazzuk az Euler-Lagrange-féle mozgásegyenletet:
\begin{equation}
    \frac{\partial \mathcal{L}}{\partial x} - \frac{d}{dt} \frac{\partial \mathcal{L}}{\partial \dot{x}} = 0
\end{equation}
\begin{equation}
    \frac{\partial \mathcal{L}}{\partial y} - \frac{d}{dt} \frac{\partial \mathcal{L}}{\partial \dot{y}} = 0
\end{equation}
Számoljuk ki az egyes tagokat:
\begin{equation}
    \frac{\partial \mathcal{L}}{\partial x} = m\omega^2 x + \frac{m\omega}{2} \dot{y} - \frac{\partial U}{\partial x}
\end{equation}
\begin{equation}
    \frac{\partial \mathcal{L}}{\partial y} = m\omega^2 y - \frac{m\omega}{2} \dot{x} - \frac{\partial U}{\partial y}
\end{equation}
\begin{equation}
    \frac{\partial \mathcal{L}}{\partial \dot{x}} = m \dot{x} - \frac{m\omega}{2} y
\end{equation}
\begin{equation}
    \frac{\partial \mathcal{L}}{\partial \dot{y}} = m \dot{y} + \frac{m\omega}{2} x
\end{equation}
\begin{equation}
    \frac{d}{dt} \frac{\partial \mathcal{L}}{\partial \dot{x}} = m \ddot{x} - \frac{m\omega}{2} \dot{y}
\end{equation}
\begin{equation}
    \frac{d}{dt} \frac{\partial \mathcal{L}}{\partial \dot{y}} = m \ddot{y} + \frac{m\omega}{2} \dot{x}
\end{equation}
Most helyettesítsük be ezeket az egyenletbe:
\begin{equation}
    m\omega^2 x + \frac{m\omega}{2} \dot{y} - \frac{\partial U}{\partial x} - (m \ddot{x} - \frac{m\omega}{2} \dot{y}) = 0
\end{equation}
\begin{equation}
    m\omega^2 y - \frac{m\omega}{2} \dot{x} - \frac{\partial U}{\partial y} - (m \ddot{y} + \frac{m\omega}{2} \dot{x}) = 0
\end{equation}
Ezt átrendezve megkapjuk a mozgásegyenleteket:
\begin{equation}
    m \ddot{x} = - \frac{\partial U}{\partial x} + m\omega^2 x + m\omega \dot{y}
\end{equation}
\begin{equation}
    m \ddot{y} = - \frac{\partial U}{\partial y} + m\omega^2 y - m\omega \dot{x}
\end{equation}
Ezek az egyenletek azt mondják ki, hogy a test gyorsulása arányos a rá ható erővel, amelyet a potenciális energia gradiensével adhatunk meg, valamint a centrifugális és Coriolis erőkkel.
Tehát a Hamilton-elvből kiindulva sikerült levezetnünk a mozgásegyenleteket egy forgó koordináta rendszerben.

\subsubsection{Lagrange-féle elsőfajú és másodfajú egyenletek}
A 18. század második felében Jean le Rond d'Alembert francia matematikus és fizikus is vizsgálta ennek az új megközelítésnek a lehetőségeit és felismerte, hogy a virtuális munka elvét ki lehet terjeszteni dinamikai rendszerekre is.
Ő azzal az ötlettel állt elő, hogy egy dinamikai rendszerben akkor beszélhetünk egyensúlyi állapotról, ha a rendszerre ható erők és a tehetetlenségi erők egyensúlyban vannak. Ezért bevezette a tehetetlenségi erők fogalmát, amelyeket a rendszer gyorsulása okoz:
\begin{equation}
    \bold{F}_{\text{tehetetlenség}} = -m \bold{a}
\end{equation}
Ahol $m$ a test tömege, és $\bold{a}$ a test gyorsulása. Ekkor a virtuális munka elve a következőképpen írható fel dinamikai rendszerekre:
\begin{equation}
    \sum_i (\bold{F}_i + \bold{F}_{\text{tehetetlenség}}) \cdot \delta \bold{r}_i = 0
\end{equation}
\begin{equation}
    \sum_i (\bold{F}_i - m_i \bold{a}_i) \cdot \delta \bold{r}_i = 0
\end{equation}
\begin{equation}
    \sum_i (\bold{F}_i - \bold{\dot{p}}_i) \cdot \delta \bold{r}_i = 0
\end{equation}
Ez onnan ered, hogy a tehetetlenségi erőket D'Alembert azért vezette be, hogy a dinamikai rendszert "egyensúlyi rendszerként" kezelhesse, ahol a virtuális munka elve alkalmazható. Ezt az elvet D'Alembert elvének nevezzük.
Ebbe a formalizmusba beilleszthetjük a kényszerfeltételeket is a Euler-Lagrange egyenlet alapján, hasonlóan ahogy azt a virtuális munka elvében tettük:
\begin{equation}
    \sum_i (\bold{F}_i - \bold{\dot{p}}_i) \cdot \delta \bold{r}_i + \sum_j \lambda_j \sum_i \boldsymbol{\nabla}_i \phi_j(\bold{r}_i) \cdot \delta \bold{r}_i = 0  
\end{equation}
\begin{equation}
    \bold{F}_i - \bold{\dot{p}}_i + \sum_j \lambda_j \boldsymbol{\nabla}_i \phi_j(\bold{r}_i) = 0
\end{equation}
\begin{equation}
    \bold{\dot{p}}_i = \bold{F}_i + \sum_j \lambda_j \boldsymbol{\nabla}_i \phi_j(\bold{r}_i)
\end{equation}
\begin{equation}
    m_i \bold{\ddot{r}}_i = \bold{F}_i + \sum_j \lambda_j \boldsymbol{\nabla}_i \phi_j(\bold{r}_i)
\end{equation}
Ez az egyenlet a rendszer mozgásegyenleteit írja le kényszerfeltételek mellett, ahol a $\lambda_j$ multiplikátorok a kényszererőket reprezentálják. Ezt az egyenletet Lagrange-féle elsőfajú egyenletnek nevezzük.

\subsubsection*{Lagrange-féle másodfajú egyenletek}
A kényszereket nem csak Lagrange multiplikátorokkal lehet kezelni, hanem úgy is, hogy a kényszerfeltételeknek megfelelően választjuk meg a koordinátákat. Ezzel a módszerrel a kényszerfeltételek automatikusan teljesülnek, és nem kell külön kényszererőket bevezetni.
Irjuk át a D'Alembert elvét általánosított koordinátákra $q_k$:
\begin{equation}
    \sum_i (\bold{F}_i - \bold{\dot{p}}_i) \cdot \delta \bold{r}_i = 0
\end{equation}
\begin{equation}
    \sum_i (\bold{F}_i - \bold{\dot{p}}_i) \cdot \sum_k \frac{\partial \bold{r}_i}{\partial q_k} \delta q_k = 0
\end{equation}
\begin{equation}
    \sum_k \left( \sum_i (\bold{F}_i - \bold{\dot{p}}_i) \cdot \frac{\partial \bold{r}_i}{\partial q_k} \right) \delta q_k = 0
\end{equation}
Mivel a virtuális elmozdulások tetszőlegesek, ezért a fenti egyenlet csak akkor teljesülhet, ha az egyes tagok külön-külön is nulla:
\begin{equation}
    \sum_i (\bold{F}_i - \bold{\dot{p}}_i) \cdot \frac{\partial \bold{r}_i}{\partial q_k} = 0
\end{equation}
Most bontsuk fel az egyes tagokat:
\begin{equation}
    \sum_i \bold{F}_i \cdot \frac{\partial \bold{r}_i}{\partial q_k} - \sum_i \bold{\dot{p}}_i \cdot \frac{\partial \bold{r}_i}{\partial q_k} = 0
\end{equation}
Ahol az első tag az általánosított erő $Q_k$:
\begin{equation}
    Q_k = \sum_i \bold{F}_i \cdot \frac{\partial \bold{r}_i}{\partial q_k}
\end{equation}
Ezt visszahelyettesítve a D'Alembert elvébe:
\begin{equation}
    \left( Q_k - \sum_i \bold{\dot{p}}_i \cdot \frac{\partial \bold{r}_i}{\partial q_k} \right) = 0
\end{equation}
Itt vizsgáljuk meg a második tagot kicsit részletesebben. Felfedezhetjük, hogy ez a kifejezés a kinetikus energia $K$ deriváltjával kapcsolatos:
\begin{equation}
    K = \sum_i \frac{1}{2} m_i \dot{\bold{r}}_i^2
\end{equation}
Ha ezt deviáljuk $q_k$ és $\dot{q}_k$ szerint, akkor a következőt kapjuk:
\begin{equation}
    \frac{\partial K}{\partial q_k} = \sum_i m_i \dot{\bold{r}}_i \cdot \frac{\partial \dot{\bold{r}}_i}{\partial q_k}
\end{equation}
\begin{equation}
    \frac{\partial K}{\partial \dot{q}_k} = \sum_i m_i \dot{\bold{r}}_i \cdot \frac{\partial \bold{\dot{r}}_i}{\partial \dot{q_k}}
\end{equation}
Ha pedig kivonjuk egymásból a kettőt, akkor a következőt kapjuk:
\begin{equation}
    \frac{d}{dt} \frac{\partial K}{\partial \dot{q}_k} - \frac{\partial K}{\partial q_k} = \sum_i m_i \bold{\ddot{r}}_i \cdot \frac{\partial \bold{r}_i}{\partial q_k} = \sum_i \bold{\dot{p}}_i \cdot \frac{\partial \bold{r}_i}{\partial q_k}
\end{equation}
Ez pedig pont a második tag a D'Alembert elvében. Ezt visszahelyettesítve:
\begin{equation}
    Q_k - \left( \frac{d}{dt} \frac{\partial K}{\partial \dot{q}_k} - \frac{\partial K}{\partial q_k} \right) = 0
\end{equation}
\begin{equation}
    \frac{d}{dt} \frac{\partial K}{\partial \dot{q}_k} - \frac{\partial K}{\partial q_k} = Q_k
\end{equation}
Ha pedig ide behelyettesítjük a konzervatív erőt, mint a potenciális energia gradiensét, akkor a következőt kapjuk:
\begin{equation}
    Q_k = - \frac{\partial U}{\partial q_k}
\end{equation}
\begin{equation}
    \frac{d}{dt} \frac{\partial K}{\partial \dot{q}_k} - \frac{\partial K}{\partial q_k} = - \frac{\partial U}{\partial q_k}
\end{equation}
\begin{equation}
    \frac{d}{dt} \frac{\partial (K - U)}{\partial \dot{q}_k} - \frac{\partial (K - U)}{\partial q_k} = 0
\end{equation}
A Kinematikus energia és a potenciális energia különbsége pedig pont a Lagrange-függvény:
\begin{equation}
    \frac{d}{dt} \frac{\partial \mathcal{L}}{\partial \dot{q}_k} - \frac{\partial \mathcal{L}}{\partial q_k} = 0
\end{equation}
Ez pedig pont az Euler-Lagrange-féle mozgásegyenlet, amely a Hamilton-elvből is levezethető. Ezt az egyenletet Lagrange-féle másodfajú egyenletnek nevezzük dinamikai rendszerek esetén.

\subsection{Szimmetriák és megmaradási tételek}
A legkisebb hatás elvének (pontosabb stacionárius hatás elve) van egy nagyon meglepő következménye, ami a Newtoni mechanikán is túlmutató eredményt hozott, és a természet egy alapvető tulajdonságát tárta fel.
Ha ugyanis a kiválasztott pályának csak annyi a feltétele, hogy a hatás stacionárius legyen, akkor bármilyen olyan transzformáció, amely a hatást nem változtatja meg, szintén egy megengedett pályát eredményez.
Például, toljuk el a kiválasztott koordináta rendszerünket egy kicsit térben:
\begin{equation}
    \bold{r} \to \bold{r} + \delta \bold{a}
\end{equation}
Ha erre felírjuk a Lagrange-függvényt, akkor azt látjuk, hogy a Lagrange-függvény nem változik meg, hiszen a kinetikus energia csak a sebességtől függ, a potenciális energia pedig csak a relatív távolságoktól.
\begin{equation}
    \mathcal{L} = K - U = \frac{1}{2}m \dot{\bold{r}}^2 - U(\bold{r}_i - \bold{r}_j)
\end{equation}
\begin{equation}
    \mathcal{L} = K - U = \frac{1}{2}m \dot{(\bold{r} + \delta \bold{a})}^2 - U((\bold{r}_i + \delta \bold{a}) - (\bold{r}_j + \delta \bold{a})) = \frac{1}{2}m \dot{\bold{r}}^2 - U(\bold{r}_i - \bold{r}_j)
\end{equation}
Tehát a Lagrange-függvény variációja nulla:
\begin{equation}
    \delta \mathcal{L} = 0
\end{equation}
Ez pedig azt jelenti, hogy a hatás sem változik meg:
\begin{equation}
    \delta S = \delta \int \mathcal{L} \, dt = 0
\end{equation}
Ez azt jelenti, hogy a rendszer "szimmetrikus" a térbeli eltolásva nézve, nem veszi észre, ha a koordináta-rendszer pontjait eltoljuk. 

\subsubsection*{Példa: forgatás}
Vegyünk egy rendszert, melyet elforgatunk $\theta$ szöggel a középpontja körül. A koordináták a következőképpen változnak:
\begin{equation}
    x' = x \cos(\theta) + y \sin(\theta)
\end{equation}
\begin{equation}
    y' = x \sin(\theta) + y \cos(\theta)
\end{equation}
A Lagrange-függvény pedig a következőképpen alakul:
\begin{equation}
    \mathcal{L} = K - U = \frac{1}{2}m(\dot{x}^2 + \dot{y}^2) - U(\sqrt{x^2 + y^2})
\end{equation}
Ha a $\theta = \delta$ szög kicsi, akkor felírhatjuk a következőképpen:
\begin{equation}
        cos(\theta) \approx 1
\end{equation}
\begin{equation}
        sin(\theta) \approx \delta
\end{equation}
Akkor a koordiniáták a következőképpen alakulnak:
\begin{equation}
    x' = x + y \delta
\end{equation}
\begin{equation}
    y' = -x \delta + y
\end{equation}
Ezt rendezve:
\begin{equation}
    \delta x = x' - x = y \delta
\end{equation}
\begin{equation}
    \delta y = y' - y = - x \delta
\end{equation}
\begin{equation}
    \delta \dot{x} = \dot{x'} - \dot{x} = \dot{y} \delta
\end{equation}
\begin{equation}
    \delta \dot{y} = \dot{y'} - \dot{y} = - \dot{x} \delta
\end{equation}
Innen a potenciális energia változása:
\begin{equation}
    \delta \sqrt{x^2 + y^2} = \frac{1}{2\sqrt{x^2 + y^2}} 2(x \delta x + y \delta y) = \frac{1}{\sqrt{x^2 + y^2}} (xy \delta - yx \delta) = 0
\end{equation}
\begin{equation}
    \rightarrow \delta U = 0
\end{equation}
A kinetikus energia változása pedig:
\begin{equation}
    \delta K = \frac{1}{2}m(2\dot{x} \delta \dot{x} + 2\dot{y} \delta \dot{y}) = m(\dot{x} \dot{y} \delta - \dot{y} \dot{x} \delta) = 0
\end{equation}
\begin{equation}
    \rightarrow \delta K = 0
\end{equation}
Tehát a Lagrange-függvény változása:
\begin{equation}
    \delta \mathcal{L} = \delta K - \delta U = 0
\end{equation}
Ez pedig azt jelenti, hogy a hatás sem változik meg:
\begin{equation}
    \delta S = \delta \int \mathcal{L} \, dt = 0
\end{equation}
Tehát a rendszer szimmetrikus forgatásra.

\subsection{Szimmetriák általánosan}
Térjünk vissza a tanszormációkra egy kicsit. Vegyünk egy transzlációs eltolást egy részecske helyzetében. Ekkor kétféle képpen értelmezhetjük a transzformációt:
\begin{itemize}
    \item Az első megközelítés szerint az egész koordináta rendszert eltoljuk, tehát a részecske helyzete változatlan marad, de a koordináták új értéket kapnak, ezt passzív transzformációnak nevezzük.
    \item A második megközelítés szerint a koordináta rendszert változatlanul hagyjuk, de a részecske helyzete változik, ezt aktív transzformációnak nevezzük.
\end{itemize}
Mindkét megközelítés helyes, és a fizikai eredményeknek függetlennek kell lenniük attól, hogy melyik megközelítést alkalmazzuk. Viszont a két lehetőség között van egy fontos különbség: az aktív transzformáció
esetén a részecske potenciális energiája is megváltozhat, hiszen a potenciális energia a részecske helyzetétől függ. Ezzel szemben a passzív transzformáció esetén a potenciális energia nem változik, hiszen a részecske helyzete változatlan marad,
csak a koordinátákat paraméterezzük át. Ennélfogva a passzív transzormáció jelen esetben nem hordoz túl sok érdekességet magában, ezért vizsgáljuk meg az aktív transzformációt egy kicsit részletesebben.
Legyen egy részecske, amelyet a konfiguációs térben egy $q_i$ általános koordináta ír le, és mozgassuk el egy kicsit $\delta q_i$-val:
\begin{equation}
    q_k \to q_k + \delta q_k
\end{equation}
\begin{equation}
    \dot{q}_k \to \dot{q}_k + \delta \dot{q}_k
\end{equation}
Akkor a Lagrange-függvény változása:
\begin{equation}
    \delta \mathcal{L} = \sum_k \left( \frac{\partial \mathcal{L}}{\partial q_k} \delta q_k + \frac{\partial \mathcal{L}}{\partial \dot{q}_k} \delta \dot{q}_k \right)
\end{equation}
Ha a rendszer szimmetrikus a transzformációra, akkor a Lagrange-függvény változása nulla:
\begin{equation}
    \delta \mathcal{L} = 0
\end{equation}
Ezt behelyettesítve:
\begin{equation}
    \sum_k \left( \frac{\partial \mathcal{L}}{\partial q_k} \delta q_k + \frac{\partial \mathcal{L}}{\partial \dot{q}_k} \delta \dot{q}_k \right) = 0
\end{equation}
Most alkalmazzuk az Euler-Lagrange-féle mozgásegyenletet:
\begin{equation}
    \frac{\partial \mathcal{L}}{\partial q_k} = \frac{d}{dt} \frac{\partial \mathcal{L}}{\partial \dot{q}_k}
\end{equation}
Ezt behelyettesítve:
\begin{equation}
    \sum_k \left( \frac{d}{dt} \frac{\partial \mathcal{L}}{\partial \dot{q}_k} \delta q_k + \frac{\partial \mathcal{L}}{\partial \dot{q}_k} \delta \dot{q}_k \right) = 0
\end{equation}
Most alkalmazzuk a természettudományokban gyakran használt "természetes" sorrendet, ahol a deriváltakat előre hozzuk:
\begin{equation}
    \sum_k \left( \frac{d}{dt} \left( \frac{\partial \mathcal{L}}{\partial \dot{q}_k} \delta q_k \right) - \frac{\partial \mathcal{L}}{\partial \dot{q}_k} \frac{d}{dt} \delta q_k + \frac{\partial \mathcal{L}}{\partial \dot{q}_k} \delta \dot{q}_k \right) = 0
\end{equation}
\begin{equation}
    \sum_k \left( \frac{d}{dt} \left( \frac{\partial \mathcal{L}}{\partial \dot{q}_k} \delta q_k \right) + \frac{\partial \mathcal{L}}{\partial \dot{q}_k} (\delta \dot{q}_k - \frac{d}{dt} \delta q_k) \right) = 0
\end{equation}
Ahol a második tag nulla, hiszen a variáció és a deriválás felcserélhető:
\begin{equation}
    \delta \dot{q}_k = \frac{d}{dt} \delta q_k
\end{equation}
Tehát a következőt kapjuk:
\begin{equation}
    \sum_k \frac{d}{dt} \left( \frac{\partial \mathcal{L}}{\partial \dot{q}_k} \delta q_k \right) = 0
\end{equation}
\begin{equation}
    \frac{d}{dt} \left( \sum_k \frac{\partial \mathcal{L}}{\partial \dot{q}_k} \delta q_k \right) = 0
\end{equation}
Ez azt jelenti, hogy a következő mennyiség megmarad:
\begin{equation}
    \sum_k \frac{\partial \mathcal{L}}{\partial \dot{q}_k} \delta q_k = \text{konstans}
\end{equation}
Ez a Noether-tétel, amely kimondja, hogy minden szimmetria esetén megmaradási törvény létezik. Például, ha a rendszer szimmetrikus térbeli eltolásra, akkor a lendület megmarad, ha szimmetrikus időbeli eltolásra, akkor az energia megmarad,
és ha szimmetrikus forgatásra, akkor a perdület megmarad. Ez egy borzasztó fontos felismerés, hisz ez a tétel köti össze a szimmetriákat a megmaradási törvényekkel, amelyek a fizika alapvető törvényei. Tehát ha egy tetszőleges szimmetriát találunk egy fizikai rendszerben,
akkor biztosak lehetünk benne, hogy létezik egy megmaradó mennyiség is.

\subsubsection*{Példa: transzlációs szimmetria}
Legyen egy rendszer ahol van $q_1$ és $q_2$ részecske, és ezeket kis $\delta$-val eltoljuk:
\begin{equation}
    \delta q_1 = \delta
\end{equation}
\begin{equation}
    \delta q_2 = \delta
\end{equation}
\begin{equation}
    \delta q_i = f_i(q) \delta
\end{equation}
Akkor a Noether-tétel szerint a következő mennyiség megmarad:
\begin{equation}
    \sum_k \frac{\partial \mathcal{L}}{\partial \dot{q}_k} \delta q_k = \sum_k p_k f_k(q) \delta = \text{konstans}
\end{equation}
\begin{equation}
    \sum_k p_k f_k(q) = p_1 + p_2 = \text{konstans}
\end{equation}
Ez a mennyiség a rendszer lendülete, amely megmarad a térbeli eltolás szimmetriája miatt.

Eltérő eltolásokra:
\begin{equation}
    \delta q_1 = b\delta
\end{equation}
\begin{equation}
    \delta q_2 = - a\delta
\end{equation}
\begin{equation}
    \delta q_i = f_i(q) \delta
\end{equation}
Akkor a Noether-tétel szerint a következő mennyiség megmarad:
\begin{equation}
    \sum_k \frac{\partial \mathcal{L}}{\partial \dot{q}_k} \delta q_k = \sum_k p_k f_k(q) \delta = \text{konstans}
\end{equation}
\begin{equation}
    \sum_k p_k f_k(q) = b p_1 - a p_2 = \text{konstans}
\end{equation}
Ez a mennyiség a rendszer impulzusmomentuma, amely megmarad a térbeli eltolás szimmetriája miatt.

Forgásra pedig:
\begin{equation}
    f_x = y
\end{equation}
\begin{equation}
    f_y = -x
\end{equation}
\begin{equation}
    p_x y - p_y x = L_z = \text{konstans}
\end{equation}
Ez a mennyiség a rendszer perdületmomentuma, amely megmarad a forgás szimmetriája miatt.

\subsubsection*{Példa: Harmonikus oszcillátor}
Legyen egy harmonikus oszcillátor, amelynek a Lagrange-függvénye a következő:
\begin{equation}
    \mathcal{L} = \frac{1}{2}m \dot{q}^2 - \frac{1}{2}k q^2
\end{equation}
Ez a rendszer időbeli eltolásra szimmetrikus, hiszen a Lagrange-függvény nem függ kifejezetten az időtől. Legyenek $m = 1, k=1$ az egyszerűség kedvéért. Tehát alkalmazhatjuk a Noether-tételt az időbeli eltolásra:
\begin{equation}
    \delta q = \dot{q} \delta t
\end{equation}
Akkor a Noether-tétel szerint a következő mennyiség megmarad:
\begin{equation}
    \sum_k \frac{\partial \mathcal{L}}{\partial \dot{q}_k} \delta q_k = \frac{\partial \mathcal{L}}{\partial \dot{q}} \delta q = p \dot{q} \delta t = \text{konstans}
\end{equation}
\begin{equation}
    p \dot{q} = m \dot{q}^2 = 2K = \text{konstans}
\end{equation}
Ez a mennyiség a rendszer kinetikus energiája, amely megmarad az időbeli eltolás szimmetriája miatt. Mivel a teljes energia $E = K + U$ is megmarad, ezért a potenciális energia is megmarad:
\begin{equation}
    U = \frac{1}{2}k q^2 = \text{konstans}
\end{equation}

\subsection{Hamilton függvény, kanonikus egyenletek, kanonikus transzformációk}
Láthattuk, hogy a Lagrange formalizmus központi eleme a Lagrange-függvény, amit eredetileg egy olyan függvényként definiáltunk, melynek az integrálja két pont között, a hatás. De hogyan is lehet értelmezni a Lagrange-függvényt?
Egyik intuitív gondolat lehet az, hogy esetleg a Lagrange-függvény a rendszer energiáját írja le. De ahhoz, hogy ez igaz legyen, egy magmaradó mennyiségnek kell lennie időben, úgyhogy vizsgáljuk meg, hogy ez helytálló-e.
Vegyünk egy általános Lagrange-függvényt, amely nem függ kifejezetten az időtől (különben nem lenne megmaradó mennyiség) és nézzük meg a deriváltját idő szerint:
\begin{equation}
    \frac{d\mathcal{L}}{dt} = \sum_k \left( \frac{\partial \mathcal{L}}{\partial q_k} \dot{q}_k + \frac{\partial \mathcal{L}}{\partial \dot{q}_k} \ddot{q}_k \right)
\end{equation}
Most alkalmazzuk az Euler-Lagrange-féle mozgásegyenletet:
\begin{equation}
    \frac{\partial \mathcal{L}}{\partial q_k} = \frac{d}{dt} \frac{\partial \mathcal{L}}{\partial \dot{q}_k}
\end{equation}
Ezt behelyettesítve:
\begin{equation}
    \frac{d\mathcal{L}}{dt} = \sum_k \left( \frac{d}{dt} \frac{\partial \mathcal{L}}{\partial \dot{q}_k} \dot{q}_k + \frac{\partial \mathcal{L}}{\partial \dot{q}_k} \ddot{q}_k \right)
\end{equation}
hozzuk be a kanonikus impulzust $ p_k = m\dot{q}_k = \frac{\partial \mathcal{L}}{\partial \dot{q}_k} $:
\begin{equation}
    \frac{d\mathcal{L}}{dt} = \sum_k \left( \dot{p}_k \dot{q}_k + p_k \ddot{q}_k \right)
\end{equation}
Ezt átrendezve:
\begin{equation}
    \frac{d\mathcal{L}}{dt} = \sum_k \frac{d}{dt} (p_k \dot{q}_k)
\end{equation}
\begin{equation}
    \frac{d\mathcal{L}}{dt} = \frac{d}{dt} \left( \sum_k p_k \dot{q}_k \right)
\end{equation}
Tehát a Lagrange-függvény nem egy megmaradó mennyiség időben, hiszen a deriváltja nem nulla, hanem a kanonikus impulzus és sebesség szorzatának az idő szerinti deriváltja. De ha ezt átrendezzük, akkor a következőt kapjuk:
\begin{equation}
    \frac{d}{dt} \left( \sum_k p_k \dot{q}_k - \mathcal{L} \right) = 0
\end{equation}
\begin{equation}
    \sum_k p_k \dot{q}_k - \mathcal{L} = \mathcal{H} = \text{konstans}
\end{equation}
Ez a mennyiség pedig már egy megmaradó mennyiség, amelyet Hamilton-függvénynek nevezünk. A Hamilton-függvény tehát a kanonikus impulzus és sebesség szorzatának és a Lagrange-függvény különbsége.
Most nézzük meg, hogy a Hamilton-függvény valóban az energia-e. Vegyünk egy egyszerű rendszert, ahol a Lagrange-függvény a következő:
\begin{equation}
    \mathcal{L} = K - U = \frac{1}{2}m \dot{q}^2 - U(q)
\end{equation}
Akkor a kanonikus impulzus:
\begin{equation}
    p = \frac{\partial \mathcal{L}}{\partial \dot{q}} = m \dot{q}
\end{equation}
A Hamilton-függvény pedig:
\begin{equation}
    \mathcal{H} = p \dot{q} - \mathcal{L} = m \dot{q}^2 - \left( \frac{1}{2}m \dot{q}^2 - U(q) \right)
\end{equation}
\begin{equation}
    \mathcal{H} = \frac{1}{2}m \dot{q}^2 + U(q) = K + U = E
\end{equation}
Tehát a Hamilton-függvény valóban az energia, amely megmarad időben, ha a Lagrange-függvény nem függ kifejezetten az időtől.

Általános esetben, ha a Lagrange-függvény expliciten is függ az időtől, akkor a következőképpen alakul a Hamilton-függvény:
\begin{equation}
    \frac{d\mathcal{L}}{dt} = \sum_k \left( \frac{\partial \mathcal{L}}{\partial q_k} \dot{q}_k + \frac{\partial \mathcal{L}}{\partial \dot{q}_k} \ddot{q}_k \right) + \frac{\partial \mathcal{L}}{\partial t}
\end{equation}
\begin{equation}
    \frac{d\mathcal{L}}{dt} = \sum_k \left( \dot{p}_k \dot{q}_k + p_k \ddot{q}_k \right) + \frac{\partial \mathcal{L}}{\partial t}
\end{equation}
\begin{equation}
    \frac{d\mathcal{L}}{dt} = \sum_k \frac{d}{dt} (p_k \dot{q}_k) + \frac{\partial \mathcal{L}}{\partial t}
\end{equation}
\begin{equation}
    \frac{d}{dt} \left( \sum_k p_k \dot{q}_k - \mathcal{L} \right) = - \frac{\partial \mathcal{L}}{\partial t}
\end{equation}
\begin{equation}
    \frac{d\mathcal{H}}{dt} = - \frac{\partial \mathcal{L}}{\partial t}
\end{equation}
Tehát ha a Lagrange-függvény kifejezetten függ az időtől, akkor a Hamilton-függvény nem marad meg időben, tehát az energia nem lesz megmaradó mennyiség!

\subsubsection{Kanonikus egyenletek}
Most, hogy megvan a Hamilton-függvény, nézzük meg, hogy mi történik ha a Hamilton-függvénybe hozunk be egy kis változást a kanonikus impulzusok és koordináták szerint:
\begin{equation}
    \delta F(q, p) = \frac{\partial F}{\partial q} \delta q + \frac{\partial F}{\partial p} \delta p \quad \rightarrow \quad \text{Ahol F egy tetszőleges függvény}
\end{equation}
\begin{equation}
    \mathcal{H} = \sum_k p_k \dot{q}_k - \mathcal{L}
\end{equation}
\begin{equation}
    \delta \mathcal{H} = \sum_k \left( \dot{q}_k \delta p_k + p_k \delta \dot{q}_k - \frac{\partial \mathcal{L}}{\partial q_k} \delta q_k - \frac{\partial \mathcal{L}}{\partial \dot{q}_k} \delta \dot{q}_k \right)
\end{equation}
Ahol az Euler-Lagrange-féle mozgásegyenletet alkalmazva:
\begin{equation}
    \frac{\partial \mathcal{L}}{\partial q_k} = \frac{d}{dt} \frac{\partial \mathcal{L}}{\partial \dot{q}_k} = \dot{p}_k
\end{equation}
\begin{equation}
    \delta \mathcal{H} = \sum_k \left( \dot{q}_k \delta p_k + p_k \delta \dot{q}_k - \dot{p}_k \delta q_k - p_k \delta \dot{q}_k \right)
\end{equation}
\begin{equation}
    \delta \mathcal{H} = \sum_k \left( \dot{q}_k \delta p_k - \dot{p}_k \delta q_k \right)
\end{equation}
Ebből következik, hogy a Hamilton-függvény parciális deriváltjai a kanonikus koordináták és impulzusok szerint a következők:
\begin{equation}
    \frac{\partial \mathcal{H}}{\partial p_k} = \dot{q}_k \quad \quad \frac{\partial \mathcal{H}}{\partial q_k} = -\dot{p}_k
\end{equation}
Ezek a Hamilton-féle kanonikus egyenletek, amelyek a Lagrange-féle mozgásegyenletek alternatív megfogalmazásai.

Tehát összefoglalva:
\begin{equation}
    \mathcal{L}(q, \dot{q}, t) \rightarrow
\end{equation}
\begin{equation}
    p = \frac{\partial \mathcal{L}}{\partial \dot{q}} \quad \quad \dot{p} = \frac{\partial \mathcal{L}}{\partial q}
\end{equation}
\begin{equation}
    \frac{d\mathcal{L}}{dt} = \sum_k \left( \frac{\partial \mathcal{L}}{\partial q_k} \dot{q}_k + \frac{\partial \mathcal{L}}{\partial \dot{q}_k} \ddot{q}_k \right) + \frac{\partial \mathcal{L}}{\partial t}
\end{equation}
\noindent\hfil\rule{0.5\textwidth}{.4pt}\hfil
\begin{equation}
    \mathcal{H}(q, p, t) = \sum_k p_k \dot{q}_k - \mathcal{L}
\end{equation}
\begin{equation}
    \frac{d\mathcal{H}}{dt} = - \frac{\partial \mathcal{L}}{\partial t}
\end{equation}
\begin{equation}
    \frac{\partial \mathcal{H}}{\partial p_k} = \dot{q}_k \quad \quad \frac{\partial \mathcal{H}}{\partial q_k} = -\dot{p}_k
\end{equation}
\noindent\hfil\rule{0.5\textwidth}{.4pt}\hfil
\begin{equation}
    \text{Kanonikus egyenletek:}
\end{equation}
\begin{equation}
    \dot{q}_k = \frac{\partial \mathcal{H}}{\partial p_k} \quad \quad \dot{p}_k = -\frac{\partial \mathcal{H}}{\partial q_k}
\end{equation}

\subsubsection{Kanonikus transzformációk}
A Hamilton-függvény és a kanonikus egyenletek bevezetése után vizsgáljuk meg, hogy milyen transzformációkat alkalmazhatunk a kanonikus koordinátákra $(q_k, p_k)$ úgy, hogy a kanonikus egyenletek formája megmaradjon.
Ezeket a transzformációkat kanonikus transzformációknak nevezzük. Legyen egy kanonikus transzformáció, amely a következőképpen néz ki:
\begin{equation}
    (q_k, p_k) \to (Q_k, P_k)
\end{equation}
A kanonikus egyenletek formája a következő:
\begin{equation}
    \dot{q}_k = \frac{\partial \mathcal{H}}{\partial p_k} \quad \quad \dot{p}_k = -\frac{\partial \mathcal{H}}{\partial q_k}
\end{equation}
\begin{equation}
    \dot{Q}_k = \frac{\partial \mathcal{K}}{\partial P_k} \quad \quad \dot{P}_k = -\frac{\partial \mathcal{K}}{\partial Q_k}
\end{equation}
Ahol $\mathcal{K}$ az új Hamilton-függvény, amely a kanonikus transzformáció után keletkezik. A kanonikus transzformáció akkor megengedett, ha a kanonikus egyenletek formája megmarad, tehát a következő feltételnek kell teljesülnie:
\begin{equation}
    \frac{\partial \mathcal{H}}{\partial p_k} = \frac{\partial \mathcal{K}}{\partial P_k} \quad \quad -\frac{\partial \mathcal{H}}{\partial q_k} = -\frac{\partial \mathcal{K}}{\partial Q_k}
\end{equation}
Ezt a feltételt úgy is felírhatjuk, hogy a következő mátrix egyenletnek kell teljesülnie:
\begin{equation}
    \begin{pmatrix}\frac{\partial Q}{\partial q} & \frac{\partial Q}{\partial p} \\ \frac{\partial P}{\partial q} & \frac{\partial P}{\partial p} \end{pmatrix}
    \begin{pmatrix} 0 & I \\ -I & 0 \end{pmatrix}
    \begin{pmatrix}\frac{\partial Q}{\partial q} & \frac{\partial Q}{\partial p} \\ \frac{\partial P}{\partial q} & \frac{\partial P}{\partial p} \end{pmatrix}^T
    = \begin{pmatrix} 0 & I \\ -I & 0 \end{pmatrix}
\end{equation}
Ahol $I$ az egységmátrix. Ez a feltétel biztosítja, hogy a kanonikus egyenletek formája megmaradjon a kanonikus transzformáció után.
Egy kanonikus transzformációt gyakran egy generátorfüggvénnyel definiálnak, amely a következőképpen néz ki:
\begin{equation}
    F(q, P, t)
\end{equation}
A generátorfüggvény segítségével a kanonikus transzformáció a következőképpen néz ki:
\begin{equation}
    p_k = \frac{\partial F}{\partial q_k} \quad \quad Q_k = \frac{\partial F}{\partial P_k}
\end{equation}
A generátorfüggvény segítségével a kanonikus transzformációk könnyen definiálhatók és kezelhetők, és számos hasznos tulajdonsággal rendelkeznek. Például, ha a generátorfüggvény nem függ kifejezetten az időtől,
akkor a kanonikus transzformáció megőrzi a Hamilton-függvény formáját, tehát az új Hamilton-függvény is megmarad időben. Ezért a kanonikus transzformációk nagyon hasznosak a Hamilton-függvények és a kanonikus egyenletek vizsgálatában,
és számos alkalmazásuk van a fizikában és a matematikában.

\subsubsection*{Példa: Harmonikus oszcillátor}
Vegyünk egy harmonikus oszcillátort, amelynek a Hamilton-függvénye a következő (teljes energia):
\begin{equation}
    \mathcal{H} = \frac{p^2}{2m} + \frac{1}{2}k q^2
\end{equation}
A kanonikus egyenletek pedig:
\begin{equation}
    \dot{q} = \frac{\partial \mathcal{H}}{\partial p} = \frac{p}{m}
\end{equation}
\begin{equation}
    \dot{p} = -\frac{\partial \mathcal{H}}{\partial q} = -k q
\end{equation}
Az impulzust mégegyszer deriválva megkaphatjuk a $F = m\ddot{q}$ alakot:
\begin{equation}
    \ddot{q} = -\frac{k}{m} q
\end{equation}
Ez pedig pont a harmonikus oszcillátor mozgásegyenlete.

\subsubsection*{Példa: Liouville-tétel}
A kanonikus transzformációk egyik fontos tulajdonsága a Liouville-tétel, amely kimondja, hogy a kanonikus transzformációk megőrzik a fázistér térfogatát. Ez azt jelenti, hogy ha egy fázistérbeli területet kanonikus transzformációnak vetünk alá,
akkor a terület mérete nem változik meg. Matematikailag ez a következőképpen írható fel:
\begin{equation}
    \int dq_1 dp_1 \ldots dq_n dp_n = \int dQ_1 dP_1 \ldots dQ_n dP_n
\end{equation}
Ez a tétel nagyon fontos a statisztikus mechanikában, ahol a fázistér térfogatának megőrzése biztosítja, hogy a rendszer állapotainak száma nem változik meg kanonikus transzformációk alatt. Ez felfogható úgy is,
hogy a fázistérben egy "folyadék" áramlik, amely nem összenyomható, tehát a térfogat nem változik meg.

\end{document}